\chapter{Econometria Espacial}

\paragraph*{Recurso experimental.} Todos os comandos \texttt{spatial\_*} deste
capítulo exigem que o Hayashi seja compilado com \texttt{--features experimental}.

\section{Por que espacial?}

Quando as observações estão localizadas no espaço ou conectadas por uma
rede, os distúrbios ou resultados podem transbordar de uma unidade para
seus vizinhos. Ignorar essa dependência produz erros padrão enviesados,
viés de coeficientes e inferência enganosa.

\section{A matriz de pesos espaciais}

Todo modelo espacial no \hayashi{} exige uma matriz de pesos $W$
padronizada por linha. Se os dados estão ordenados por unidade,
$W_{ij}$ é o peso que a unidade $j$ recebe ao prever o resultado da
unidade $i$.

\begin{lstlisting}[caption={Matriz de pesos espaciais}]
// 3 regiões com vizinhos mútuos
let W = [[0, 0.5, 0.5],
         [0.5, 0, 0.5],
         [0.5, 0.5, 0]]
\end{lstlisting}

Para grades maiores, gere $W$ a partir do DataFrame ou carregue-a de um
arquivo.

\section{SAR espacial}

O modelo Autorregressivo Espacial é

\[
y = \rho W y + X\beta + \varepsilon
\]

onde $\rho$ é o parâmetro autorregressivo espacial. A estimação usa
verossimilhança concentrada: para cada valor em uma grade de $\rho$,
transforma a variável dependente e roda OLS; depois seleciona o $\rho$
que maximiza a log-verossimilhança concentrada.

\begin{lstlisting}[caption={SAR espacial}]
let n = 100
let W = makeweights(df, "region", type="rook")  // ou lista de listas
generate df x1 = rnormal(n)
generate df y = 0.5*x1 + rnormal(n)
generate df wy = spatial_lag(W, y)
generate df y = y + 0.3*wy

let m_sar = spatial_sar(y ~ x1, df, w=W)
print(m_sar)
\end{lstlisting}

A saída reporta $\rho$, $\beta$, seus erros padrão e um teste de
verossimilhança para $\rho = 0$.

\section{SEM espacial}

O Modelo de Erro Espacial é

\[
y = X\beta + u, \qquad u = \lambda W u + \varepsilon
\]

Os erros são espacialmente correlacionados, mas o resultado em si não é
diretamente influenciado pelos resultados dos vizinhos.

\begin{lstlisting}[caption={SEM espacial}]
let m_sem = spatial_sem(y ~ x1, df, w=W)
print(m_sem)
\end{lstlisting}

\section{Durbin espacial e extensões em painel}

O Modelo de Durbin Espacial adiciona regressores defasados
espacialmente:

\[
y = \rho W y + X\beta + W X \theta + \varepsilon
\]

Ele encaixa SAR ($\theta = 0$) e o modelo de defasagem espacial
($\rho = 0$).

\begin{lstlisting}[caption={Durbin espacial}]
let m_sdm = spatial_durbin(y ~ x1, df, w=W)
print(m_sdm)
\end{lstlisting}

Para dados em painel com efeitos fixos, use as variantes de painel:

\begin{lstlisting}[caption={SAR em painel}]
let m_panelsar = spatial_panel_sar(y ~ x1, df, w=W, id="region")
print(m_panelsar)
\end{lstlisting}

\section{Interpretando efeitos diretos e indiretos}

No modelo de Durbin Espacial o coeficiente de $x$ não é o efeito total.
O efeito total na unidade $i$ inclui o impacto direto $\beta$ mais o
impacto indireto que viaja pela rede através de $\theta$ e $\rho$.

\hayashi{} reporta efeitos diretos, indiretos e totais no resumo de
\texttt{m\_sdm} quando disponíveis.

\section{Notas práticas}

\begin{itemize}
\item Sempre verifique se $W$ está padronizada por linha; $W$ não
      padronizada distorce a interpretação de $\rho$ e $\lambda$.
\item Amostras maiores são necessárias para MLE espacial do que para OLS.
\item Modelos espaciais de corte transversal assumem uma rede estável.
      Modelos espaciais em painel adicionam efeitos fixos, mas ainda
      exigem $W$ constante no tempo.
\item O modelo \texttt{spatial\_durbin\_error} combina regressores defasados
      espacialmente com um termo de erro espacial.
\end{itemize}
